% ============================================================================
% Fichier  : partie3/P03_C10_normes.tex
% Titre    : Normes et Référentiels de Sécurité
% Auteur   : MINKA MI NGUIDJOÏ Thierry Emmanuel
% Version  : 1.0 -- novembre 2026
% ============================================================================
% Sources : ISO/IEC 27001:2022 (clauses 4 à 10, annexe A de 93 mesures en
% quatre thèmes : 37 organisationnelles, 8 liées aux personnes, 14 physiques,
% 34 technologiques) ; ISO/IEC 27002:2022 ; NIST CSF 2.0 (février 2024, six
% fonctions) ; CIS Controls v8 (18 contrôles, trois groupes de mise en
% œuvre) ; ISO/IEC 27006 pour les organismes de certification.
% ============================================================================

\chapter{Normes et Référentiels de Sécurité}
\label{chap:normes}
\chaptermark{Normes et référentiels}

\epigraph{%
    « Une norme ne dit pas ce qu'il faut faire chez vous~;\\
    elle dit ce qu'il ne faut pas oublier. »%
}{Principe du cours}

\niveauChapitre{Fondamental}{Chapitres~\ref{chap:gouvernance}
et~\ref{chap:risques}.}

\begin{boxobjectifs}
\begin{itemize}
    \item Décrire la structure de l'ISO/IEC~27001 et de son annexe~A.
    \item Distinguer exigence (27001) et recommandation (27002).
    \item Présenter les six fonctions du NIST CSF~2.0 et les CIS Controls.
    \item Choisir un référentiel selon le besoin, et les combiner.
    \item Décrire le processus de certification.
\end{itemize}
\end{boxobjectifs}

\bigskip

% ============================================================================
\section*{Question directrice : par où LiZbiZ doit-elle commencer~?}
% ============================================================================

Après l'attaque, M.~ETEME reçoit trois conseils contradictoires~: « faites
certifier ISO~27001 », « appliquez le cadre du NIST », « commencez par les
contrôles CIS ». Ce chapitre explique ce que propose chaque référentiel,
pour qui il est fait, et comment les combiner sans s'y perdre.

\medskip
\textbf{Plan du chapitre.}
\begin{enumerate}
    \item Pourquoi des normes (section~\ref{sec:c10-pourquoi}).
    \item L'ISO/IEC~27001 et l'ISO/IEC~27002 (section~\ref{sec:c10-iso}).
    \item Le NIST Cybersecurity Framework~2.0 (section~\ref{sec:c10-nist}).
    \item Les CIS Controls (section~\ref{sec:c10-cis}).
    \item Choisir et combiner (section~\ref{sec:c10-choisir}).
    \item La certification (section~\ref{sec:c10-certif}).
\end{enumerate}

\separateur

% ============================================================================
\section{Pourquoi des Normes}
\label{sec:c10-pourquoi}
% ============================================================================

Un référentiel apporte une liste éprouvée de sujets à traiter, un
vocabulaire commun avec les partenaires, et une base de comparaison pour
l'audit. La loi n°~2010/012 demande aux exploitants de se doter de
« systèmes normalisés » de gestion des risques (art.~26,
Chapitre~\ref{chap:gouvernance})~: une norme reconnue est le moyen le plus
direct de le démontrer.

\begin{boxinfo}
\textbf{Norme ou référentiel~?} Une \emph{norme} est publiée par un
organisme de normalisation (ISO, AFNOR, ANOR au Cameroun)~; un
\emph{référentiel} peut venir d'une agence publique (NIST) ou d'une
organisation professionnelle (CIS). Seules certaines normes, comme
l'ISO/IEC~27001, se prêtent à une certification par un tiers.
\end{boxinfo}

% ============================================================================
\section{L'ISO/IEC 27001 et l'ISO/IEC 27002}
\label{sec:c10-iso}
% ============================================================================

\begin{boxdef}
\termecle{Système de management de la sécurité de l'information} (SMSI)~:
ensemble des politiques, procédures, ressources et activités par lesquels
une organisation gère la sécurité de son information selon un cycle
d'amélioration continue. L'ISO/IEC~27001 en fixe les
\emph{exigences}~\cite{iso27001_2022}.
\end{boxdef}

\begin{tableaucours}{L{1.6cm}L{4.4cm}L{6.2cm}}{%
    \tableauHeader{Clause} &
    \tableauHeader{Thème} &
    \tableauHeader{Exemple d'exigence}
}
4 & Contexte & Définir le périmètre du SMSI \\
5 & Leadership & Politique approuvée, rôles attribués \\
6 & Planification & Appréciation et traitement des risques, déclaration
d'applicabilité \\
7 & Support & Ressources, compétences, sensibilisation, documentation \\
8 & Fonctionnement & Mettre en œuvre le plan de traitement \\
9 & Évaluation des performances & Surveillance, audit interne, revue de
direction \\
10 & Amélioration & Non-conformités, actions correctives \\
\end{tableaucours}
\vspace{-0.8em}
\begin{center}{\small\itshape Tableau~10.1. Les clauses d'exigences de
l'ISO/IEC~27001:2022}\end{center}

L'annexe~A de l'édition 2022 compte 93~mesures, réparties en quatre
thèmes~: 37 organisationnelles, 8 liées aux personnes, 14 physiques et
34 technologiques. L'organisation retient celles qui traitent ses risques
et justifie les autres dans une \termecle{déclaration d'applicabilité}.
L'ISO/IEC~27002 détaille chaque mesure~: objectif, recommandations de mise
en œuvre~\cite{iso27002_2022}.

\begin{boxalert}
\textbf{Violation typique.} Appliquer les 93 mesures comme une liste de
cases à cocher.

\smallskip
\textbf{Règle~:} l'ISO/IEC~27001 part des risques~; les mesures en
découlent. Une mesure non pertinente s'exclut avec une justification~;
une mesure pertinente difficile ne s'exclut pas.
\end{boxalert}

% ============================================================================
\section{Le NIST Cybersecurity Framework 2.0}
\label{sec:c10-nist}
% ============================================================================

Le NIST a publié en février~2024 la version~2.0 de son
\emph{Cybersecurity Framework}~\cite{nist2024csf}, conçu pour toute
organisation. Il organise les résultats attendus en six fonctions.

\begin{tableaucours}{L{2.6cm}L{9.6cm}}{%
    \tableauHeader{Fonction} &
    \tableauHeader{Résultats attendus (exemples)}
}
Gouverner & Stratégie, rôles, politique, gestion des risques, chaîne
d'approvisionnement \\
Identifier & Inventaire des actifs, appréciation des risques,
amélioration \\
Protéger & Identités et accès, sensibilisation, sécurité des données et
des plateformes \\
Détecter & Surveillance continue, analyse des événements \\
Répondre & Gestion, analyse, communication et atténuation des incidents \\
Rétablir & Exécution du plan de reprise, communication \\
\end{tableaucours}
\vspace{-0.8em}
\begin{center}{\small\itshape Tableau~10.2. Les six fonctions du NIST CSF~2.0}\end{center}

Le cadre propose des \emph{profils}~: un profil actuel (où l'on en est) et
un profil cible (où l'on veut aller)~; l'écart entre les deux donne le
plan d'action.

% ============================================================================
\section{Les CIS Controls}
\label{sec:c10-cis}
% ============================================================================

Le Center for Internet Security publie dix-huit contrôles prioritaires,
classés en trois \emph{groupes de mise en œuvre}~: le premier (IG1)
réunit l'« hygiène » de base que toute organisation devrait appliquer, les
suivants s'adressent aux organisations plus exposées~\cite{cis2021controls}.

\begin{tableaucours}{L{1.2cm}L{6.4cm}L{4.6cm}}{%
    \tableauHeader{N°} &
    \tableauHeader{Contrôle (extrait)} &
    \tableauHeader{Opération GAROUA}
}
1 & Inventaire des actifs matériels & Machines inconnues sur le réseau \\
4 & Configuration sécurisée & Services inutiles ouverts \\
5 & Gestion des comptes & Compte partagé du prestataire \\
6 & Gestion des accès & Pas de double authentification \\
7 & Gestion continue des vulnérabilités & Correctif non installé \\
8 & Gestion des journaux & Journaux non collectés \\
11 & Récupération des données & Sauvegardes en ligne, non testées \\
15 & Gestion des prestataires & Contrat sans clause de sécurité \\
\end{tableaucours}
\vspace{-0.8em}
\begin{center}{\small\itshape Tableau~10.3. Des CIS Controls à l'Opération GAROUA}\end{center}

% ============================================================================
\section{Choisir et Combiner}
\label{sec:c10-choisir}
% ============================================================================

\begin{tableaucours}{L{2.8cm}L{3.0cm}L{3.0cm}L{3.4cm}}{%
    \tableauHeader{Référentiel} &
    \tableauHeader{Nature} &
    \tableauHeader{Certifiable} &
    \tableauHeader{Point fort}
}
ISO/IEC~27001 & Exigences de management & Oui & Reconnaissance
internationale, gouvernance \\
ISO/IEC~27002 & Recommandations & Non & Détail des mesures \\
NIST CSF~2.0 & Cadre de résultats & Non & Lisible par la direction,
profils \\
CIS Controls & Mesures priorisées & Non & Priorités concrètes, petites
structures \\
\end{tableaucours}
\vspace{-0.8em}
\begin{center}{\small\itshape Tableau~10.4. Quatre référentiels, quatre usages}\end{center}

\begin{boxscenario}{Opération GAROUA. La réponse à M. ETEME}
\textbf{Recommandation du RSSI.} Commencer par le groupe IG1 des CIS
Controls, pour corriger en trois mois les failles exploitées en mars~;
présenter au conseil la situation selon les six fonctions du NIST CSF,
plus lisibles pour lui~; construire en dix-huit mois un SMSI conforme à
l'ISO/IEC~27001, puis décider d'une certification si les clients ou les
banques l'exigent.

\medskip
\textbf{Analyse} (confiance moyenne). Les trois conseils n'étaient pas
contradictoires~: ils répondaient à trois horizons différents.
\end{boxscenario}

% ============================================================================
\section{La Certification}
\label{sec:c10-certif}
% ============================================================================

Un organisme de certification accrédité (selon l'ISO/IEC~27006) audite le
SMSI en deux étapes~: une revue documentaire, puis un audit de mise en
œuvre. Le certificat vaut trois ans, avec des audits de surveillance
annuels. Il atteste qu'un système de management fonctionne sur un
périmètre donné~; il ne garantit pas l'absence d'incident.

\begin{boxmetier}{Auditeur de sécurité}
\textbf{Lire un certificat.} L'auditeur vérifie le périmètre exact du
certificat (un seul centre de données~?), sa date, l'organisme et son
accréditation, et demande la déclaration d'applicabilité. Un certificat
dont le périmètre exclut le service audité ne prouve rien pour lui.
\end{boxmetier}

\begin{boxmetier}{RSSI}
\textbf{Faire une table de correspondance.} Le RSSI rattache chaque
mesure de son plan à la clause ou à la mesure de chaque référentiel
qu'elle satisfait~: un seul effort, plusieurs exigences couvertes.
\end{boxmetier}

\separateur

\begin{boxresume}
\textbf{Ce qu'il faut retenir de ce chapitre~:}
\begin{itemize}
    \item \textbf{L'ISO/IEC~27001} fixe les exigences d'un SMSI (clauses 4
          à~10)~; son annexe~A compte 93~mesures en quatre thèmes,
          retenues selon les risques.
    \item \textbf{L'ISO/IEC~27002} détaille les mesures~; elle ne se
          certifie pas.
    \item \textbf{Le NIST CSF~2.0} organise la sécurité en six fonctions
          (gouverner, identifier, protéger, détecter, répondre, rétablir).
    \item \textbf{Les CIS Controls} donnent 18 contrôles priorisés~; IG1 est
          l'hygiène de base.
    \item \textbf{On combine} selon l'horizon~; une certification atteste
          un système sur un périmètre, pas l'absence d'incident.
\end{itemize}
\end{boxresume}

\begin{boxinfo}
\textbf{Points de contrôle de l'élève.} Je sais~: citer les clauses 4 à
10 de l'ISO/IEC~27001~; décrire l'annexe~A~; nommer les six fonctions du
NIST CSF~2.0~; expliquer les groupes de mise en œuvre des CIS Controls~;
choisir un référentiel selon le besoin~; lire un certificat.
\end{boxinfo}

\bigskip

\begin{boxexo}[title={\ding{45}\quad Exercice~10.1. Exigence ou recommandation \niveaubase}]
Expliquez la différence entre l'ISO/IEC~27001 et l'ISO/IEC~27002, et
dites laquelle on peut certifier.
\end{boxexo}

\begin{boxexo}[title={\ding{45}\quad Exercice~10.2. Les fonctions \niveaubase}]
Classez dans les fonctions du NIST CSF~2.0~: tester une restauration~;
nommer un RSSI~; surveiller les journaux~; tenir l'inventaire~; isoler
une machine infectée~; former les utilisateurs.
\end{boxexo}

\begin{boxexo}[title={\ding{45}\quad Exercice~10.3. Annexe A \niveaubase}]
Combien de mesures compte l'annexe~A de l'ISO/IEC~27001:2022, et dans
quels thèmes~?
\end{boxexo}

\begin{boxexo}[title={\ding{45}\quad Exercice~10.4. Exclure \niveaumoyen}]
La BANQUE SANAGA veut exclure la mesure sur le développement sécurisé au
motif qu'elle « achète ses logiciels ». Est-ce acceptable~? À quelle
condition~?
\end{boxexo}

\begin{boxexo}[title={\ding{45}\quad Exercice~10.5. Profils \niveaumoyen}]
Construisez, pour la fonction « Protéger » du NIST CSF, le profil actuel
de LiZbiZ en mars et un profil cible à douze mois.
\end{boxexo}

\begin{boxexo}[title={\ding{45}\quad Exercice~10.6. IG1 \niveaumoyen}]
Pourquoi recommander les CIS Controls de niveau IG1 à une PME de
Bafoussam de 20~salariés plutôt qu'une certification ISO/IEC~27001~?
\end{boxexo}

\begin{boxexo}[title={\ding{45}\quad Exercice~10.7. Correspondance \niveaumoyen}]
Construisez une table de correspondance pour trois mesures de LiZbiZ
(double authentification, sauvegardes hors ligne, journalisation)~:
référentiel, clause ou mesure couverte.
\end{boxexo}

\begin{boxexo}[title={\ding{45}\quad Exercice~10.8. Un certificat \niveaumoyen}]
SAHEL INFOGÉRANCE présente un certificat ISO/IEC~27001 dont le périmètre
est « le centre de services de Garoua, hors prestations chez les
clients ». Qu'en déduisez-vous pour LiZbiZ~?
\end{boxexo}

\begin{boxexo}[title={\ding{45}\quad Exercice~10.9. Feuille de route \niveauexpert}]
Construisez la feuille de route de dix-huit mois de LiZbiZ vers un SMSI
conforme, par trimestre, avec les livrables de chaque clause.
\end{boxexo}

\begin{boxexo}[title={\ding{45}\quad Exercice~10.10. Choisir pour l'ANPP \niveauexpert}]
L'ANPP, agence publique, doit choisir un référentiel. Analysez ses
contraintes (moyens, obligations, attentes des partenaires) et
recommandez une combinaison.
\end{boxexo}

\bigskip

\begin{boxinfo}
\textbf{Pour aller plus loin.}
\begin{itemize}
    \item ISO/IEC~27001 et 27002:2022~\cite{iso27001_2022,iso27002_2022}.
    \item NIST, \emph{Cybersecurity Framework~2.0}~\cite{nist2024csf}.
    \item CIS Controls~v8~\cite{cis2021controls}.
\end{itemize}
\end{boxinfo}

\begin{boxinfo}
\textbf{Transition.} Les normes sont volontaires~; la loi ne l'est pas.
Le Chapitre~\ref{chap:droit} présente le droit de la sécurité au Cameroun
et en Afrique.
\end{boxinfo}
